% =====================================================================
% report.tex  -  IEEE two-column paper for AI Academy ML Final Project
% Compile with:  pdflatex report.tex && pdflatex report.tex
% NOTE: compile from the report/ folder so figure paths resolve correctly.
% =====================================================================
\documentclass[conference]{IEEEtran}
\usepackage[utf8]{inputenc}
\usepackage[T1]{fontenc}
\usepackage{amsmath,amssymb}
\usepackage{graphicx}
\usepackage{booktabs}
\usepackage{array}
\usepackage{hyperref}
\usepackage{xcolor}
\usepackage{multirow}

% ===== TITLE & AUTHORS ================================================
\title{Ensemble Methods: Boosting vs.\ Bagging\\An Empirical Study\\
{\large AI Academy Final Project Report, Summer 2026}}

\author{
\IEEEauthorblockN{Shahin Alakparov, Narmina Ibragimova, Hasan Mammadov}
\IEEEauthorblockA{AI Academy, National AI Center\\
Spring 2026\\
\{sahin.alekperov25, narmine.ibrahimova25, hasan.memmedov25\}@aiacademy.az}
}

% ======================================================================
\begin{document}
\maketitle

% -- Abstract ----------
\begin{abstract}
We present a systematic empirical comparison of two ensemble learning
paradigms-\emph{boosting} and \emph{bagging}-implemented entirely
from scratch using NumPy.  Our custom implementations include a CART
Decision Tree (with $O(N\!\log\!N)$ split search), AdaBoost (discrete
SAMME with decision stumps), and a Random Forest (bootstrap aggregation
with random feature sub-sampling and OOB scoring).  We evaluate on four
benchmark datasets spanning low-dimensional medical data, high-dimensional
image classification, large tabular census data, and a severely class-imbalanced
multiclass forest-cover problem.  Seven controlled
experiments investigate scaling, noise robustness, bias-variance
decomposition, and cluster structure.  We find that \emph{AdaBoost
minimises Bias$^2$ more aggressively (0.0231 vs.\ 0.0317 for RF) while
Random Forest achieves far lower Variance (0.0076 vs.\ 0.0131 for AdaBoost)}.
On the severely imbalanced Covertype dataset, AdaBoost's accuracy (0.6256)
masks a macro-F1 collapse to 0.3186, while Random Forest maintains both higher
accuracy (0.8501) and a far smaller accuracy-F1 gap (0.7457), exposing a failure
mode invisible to accuracy alone.
Under label corruption at $\eta{=}0.20$, Random Forest degrades only
marginally while AdaBoost suffers a significant accuracy collapse on clean,
well-structured data, confirming that re-weighting amplifies label noise
when the boundary is otherwise clean; on severely imbalanced Covertype,
injected noise instead \emph{improves} AdaBoost's accuracy by forcing it
to abandon unlearnable rare classes in favour of the dominant ones.  On
clean data both ensembles substantially outperform a single tree, with
Random Forest achieving the highest AUC-ROC on Breast Cancer (where all
figures are reported).  Our implementations
match scikit-learn baselines within 1\% absolute accuracy, confirming
correctness.  Code is available at \texttt{https://github.com/shahin1717/ml-project} (tag
\texttt{v1.0-final}).
\end{abstract}

% -- Introduction ---------
\section{Introduction}
\label{sec:intro}

\IEEEPARstart{E}{nsemble} methods are among the most successful
machine learning techniques in practice~\cite{hastie2009elements}.
By combining multiple base learners, ensembles reduce generalisation
error more effectively than any individual model.  Two dominant
strategies exist: \emph{boosting}, which builds an additive model
sequentially to reduce bias, and \emph{bagging}, which trains models
independently in parallel to reduce variance.

The central question of this project is: \textit{``Under what
conditions does boosting outperform bagging, and vice versa, and
why?''}  To answer this rigorously we implement all algorithms from
scratch, eliminating the risk of attributing differences to
implementation details rather than algorithmic properties.  We
deliberately avoid ensemble classes from scikit-learn (except as a
reference baseline) and use only NumPy for all numerical operations.

This setting is representative of real-world ML engineering: when
a team needs to port an algorithm to a new language or hardware target,
understanding the mathematical core-rather than relying on library
abstractions-is essential.

\textbf{Contributions.} (1) Fully from-scratch implementations of
CART, AdaBoost, and Random Forest with 93\% test coverage.
(2) Seven reproducible experiments with fixed random seeds on
\textbf{four} diverse datasets, including one satisfying the
specification's severe-imbalance requirement (minority class
$\le 1\%$). (3) Quantitative bias-variance decomposition
comparing boosting and bagging.  (4) Bonus implementations of
Gradient Boosting and SAMME.R with comparative analysis.

The paper is organised as follows.  Section~\ref{sec:related} reviews
the theoretical background.  Section~\ref{sec:methods} describes our
implementations.  Section~\ref{sec:experiments} details the
experimental design.  Section~\ref{sec:results} presents empirical
results with analysis.  Section~\ref{sec:discussion} synthesises
findings, and Section~\ref{sec:conclusion} concludes.

% -- Related Work ---------
\section{Related Work}
\label{sec:related}

\subsection{Decision Trees}
CART~\cite{quinlan1986induction} builds binary trees by greedily
selecting splits that minimise impurity.  The two standard impurity
measures are Gini impurity $I_G = 1 - \sum_c p_c^2$ and Shannon
entropy $I_E = -\sum_c p_c \log_2 p_c$.  Trees are prone to
overfitting on noisy data and have high variance-the primary
motivation for ensemble methods.

\subsection{AdaBoost}
Freund and Schapire~\cite{freund1997decision} showed that any
\emph{weak learner} (accuracy slightly above 50\%) can be boosted
into a strong learner.  AdaBoost maintains a distribution over
training samples, upweighting misclassified examples at each round.
The final classifier is a weighted majority vote.  Hastie \emph{et
al.}~\cite{hastie2009elements} proved that AdaBoost minimises an
exponential loss and can be interpreted as gradient boosting in
function space.

\subsection{Bagging and Random Forests}
Breiman~\cite{breiman1996bagging} introduced bootstrap aggregation
(bagging) as a general variance-reduction technique.  Random
Forests~\cite{breiman2001random} extend bagging by sub-sampling
features at each node split, which decorrelates individual trees
and dramatically reduces variance beyond what pure bagging achieves.
Out-of-bag (OOB) error provides an unbiased generalisation estimate
without requiring a separate validation set.

\subsection{Bias-Variance Trade-off}
The expected prediction error decomposes as
$\text{Err} = \text{Bias}^2 + \text{Variance} + \sigma^2_\varepsilon$.
Boosting primarily reduces bias by correcting errors iteratively;
bagging primarily reduces variance by averaging~\cite{hastie2009elements}.
Breiman~\cite{breiman1996bagging} proved that averaging reduces
variance by a factor proportional to the correlation between trees.

\subsection{Unsupervised Methods}
PCA~\cite{pearson1901lines} reduces dimensionality by projecting onto
directions of maximum variance.  K-Means~\cite{lloyd1982least}
partitions data by minimising within-cluster sum of squares.  DBSCAN~\cite{ester1996density}
identifies clusters of arbitrary shape without requiring the number of
clusters as input, labelling low-density points as noise.

% -- Methods ----------
\section{Methods}
\label{sec:methods}

\subsection{Decision Tree (CART)}
We implement binary splitting with exhaustive threshold search over
continuous features.  For each feature we sort the $N$ samples
($O(N\log N)$) and sweep all midpoints between consecutive distinct
values using \texttt{np.cumsum} to compute weighted class counts
incrementally-avoiding an inner Python loop over samples.  The
per-feature complexity is thus $O(N\log N + N)$ rather than
$O(N^2)$.  The overall split search is $O(PN\log N)$ where $P$ is
the number of candidate features.

The impurity reduction (information gain) for a candidate split is:
\begin{equation}
\Delta I = I(S) - \frac{|S_L|}{|S|}\,I(S_L) - \frac{|S_R|}{|S|}\,I(S_R)
\label{eq:gain}
\end{equation}
Stopping criteria: (i) maximum depth reached, (ii) node samples
$< \texttt{min\_samples\_split}$, (iii) pure node, (iv) all feature
vectors identical.  Feature importances are computed as normalised
mean decrease in impurity (MDI) weighted by node sample count.

\subsection{AdaBoost (Discrete SAMME)}
We implement the discrete SAMME variant~\cite{hastie2009elements},
which generalises AdaBoost to $K$-class problems.  Let
$w_i^{(1)} = 1/N$.  At each round $m$:
\begin{enumerate}
  \item Fit a stump $h_m$ on weighted data.
  \item Compute weighted error $\epsilon_m = \sum_i w_i^{(m)}\mathbf{1}[h_m(x_i)\neq y_i]$.
  \item Compute estimator weight:
    $\alpha_m = \beta\!\left(\ln\frac{1{-}\epsilon_m}{\epsilon_m} + \ln(K{-}1)\right)$
    where $\beta$ is the learning rate.
  \item Update and renormalise weights:
    $w_i^{(m+1)} \propto w_i^{(m)}\exp(\alpha_m\,\mathbf{1}[h_m(x_i)\neq y_i])$.
\end{enumerate}
Early stopping triggers when $\epsilon_m \ge (K-1)/K$ (worse than
random guessing).  Each stump receives a seed derived from the master
random state to ensure reproducibility.

\subsection{Random Forest}
Bootstrap sampling draws $N$ samples with replacement per tree;
the out-of-bag set is the complement.  At each node, $\lfloor\sqrt{P}\rfloor$
features are randomly sampled before finding the best split.  Trees
are trained in parallel using \texttt{multiprocessing.Pool}.
OOB predictions are accumulated only from trees for which a sample
was \emph{not} used during training, giving an unbiased score
$\hat{\text{acc}}_\text{OOB}$.  Feature importances average the MDI
scores across all trees.

\subsection{Unsupervised Pipeline}
\textbf{PCA} centres the data and decomposes the covariance matrix
via \texttt{np.linalg.eigh}.  \textbf{K-Means} uses Lloyd's algorithm
with $k$-distance initialisation over multiple restarts, tracking the
lowest inertia.  \textbf{DBSCAN} identifies core points (at least
\texttt{min\_samples} neighbours within $\varepsilon$) and expands
clusters via BFS.

% -- Experimental Setup --------
\section{Experimental Setup}
\label{sec:experiments}

\subsection{Datasets}
Table~\ref{tab:datasets} summarises the four datasets.  Each
presents different modelling challenges.

\begin{table}[h]
\centering
\caption{Dataset characteristics after preprocessing.}
\label{tab:datasets}
\setlength{\tabcolsep}{4pt}
\begin{tabular}{lcccl}
\toprule
\textbf{Dataset} & \textbf{N} & \textbf{P} & \textbf{Classes} & \textbf{Challenge}\\
\midrule
Breast Cancer & 569 & 30 & 2 & Small, clean\\
Adult Income & 45\,222 & 100 & 2 & Large, imbalanced\\
MNIST 3-vs-8 & 6\,000 & 784 & 2 & High-dim, visual\\
Covertype & 30\,000 & 54 & 7 & Severe imbalance ($\le$1\% minority)\\
\bottomrule
\end{tabular}
\end{table}

\textbf{Breast Cancer Wisconsin} is a small, clean binary dataset
predicting malignancy from 30 real-valued cell-nucleus measurements.
The low noise and clean boundary make it an ideal benchmark.

\textbf{Adult Income} predicts whether annual income exceeds \$50K.
The dataset is class-imbalanced (${\approx}76\%$ negative class),
contains mixed continuous and categorical features (one-hot encoded
to 100 columns after one-hot expansion of 8 categorical variables). Rows with missing values (marked \texttt{?})
were discarded via \texttt{dropna}, reducing the dataset size from 48,842 to
45,222 cleaned examples.

\textbf{MNIST 3-vs-8} is a 784-dimensional binary classification
problem.  Digits 3 and 8 share curved strokes, making them one of
the hardest MNIST binary pairs.  We subsample 6,000 images
(${\approx}3000$ per class) to satisfy the specification requirement
of $\ge 5000$ samples while keeping computation manageable.

\textbf{Covertype} predicts forest cover type from cartographic
variables across seven classes. The full dataset's rarest class
(Cottonwood/Willow) constitutes $\approx$0.5\% of samples; we
subsample 30{,}000 rows via \textbf{simple random sampling}
(not stratified) to preserve this natural rarity. We confirmed the
rarest class remains at $0.457\%$ ($137$/$30{,}000$) of our subsample, satisfying the specification's
requirement of at least one dataset with severe class imbalance
(minority class $\le$1\%). No class-rebalancing technique
(oversampling, SMOTE, class weights) was applied, since the goal of
this dataset is precisely to expose how each ensemble
\emph{natively} handles severe imbalance without correction.

\subsection{Preprocessing}
Features were standardised (zero mean, unit variance) using training
statistics only, applied to the test set identically.  No data
augmentation was applied.

\subsection{Evaluation Protocol}
Experiment~1 (correctness) uses a single fixed 80/20 train-test split.
Experiments~2-3 (scaling curves) use a single fixed split to enable
staged ensemble predictions across estimator counts.  Experiment~4
(head-to-head) uses stratified 5-fold cross-validation with
\texttt{random\_state=42}.  Experiment~5 (noise robustness) uses a
single fixed split, varying only the label corruption fraction.
We report accuracy, macro F1-score,
and AUC-ROC.  Experiment~6 (bias-variance) uses $B{=}100$ bootstrap
replicates of the Breast Cancer training set, evaluating all models
on a fixed held-out test split.  All results reported below used
100 estimators for both ensembles.

% -- Results ----------
\section{Results}
\label{sec:results}

\subsection{Experiment 1: Correctness Verification}
\label{ssec:baseline}

We compare our custom \texttt{DecisionTree} against
\texttt{sklearn.tree.DecisionTreeClassifier} (Table~\ref{tab:baseline}).
Crucially, all differences are within 1\%, satisfying the $\le 2\%$
threshold required by the specification.

\begin{table}[h]
\centering
\caption{Experiment~1: Baseline accuracy-custom vs.\ sklearn.}
\label{tab:baseline}
\setlength{\tabcolsep}{4pt}
\begin{tabular}{lcccc}
\toprule
\textbf{Dataset} & \textbf{Custom} & \textbf{Sklearn} & \textbf{Stump} & $|\Delta|$\\
\midrule
Breast Cancer & 0.9469 & 0.9469 & 0.8850 & 0.0000\\
Adult Income  & 0.8092 & 0.8105 & 0.7543 & 0.0013\\
MNIST 3-vs-8  & 0.9292 & 0.9317 & 0.8392 & 0.0025\\
Covertype     & 0.7802 & 0.7780 & 0.6388 & 0.0022\\
\bottomrule
\end{tabular}
\end{table}

\textbf{Why is Adult accuracy ``only'' 81\%?}  Adult Income is
genuinely difficult: the negative class accounts for $76\%$ of
samples, the feature space is high-dimensional after OHE (100 columns),
and the income boundary is a complex non-linear function of age, education,
occupation, and capital gains.  A single unpruned tree achieves 81\%
partly by defaulting toward the majority class.  The custom tree
matches sklearn within 0.2\% (0.8092 vs.\ 0.8105), well within the
$\le 2\%$ threshold.  The stump at 75\%
barely exceeds the na\"ive majority-class baseline ($\approx76\%$).

\textbf{Why does the stump achieve only F1=0.43 on Adult?}  A depth-1
tree can only partition the 100-dimensional space with one binary cut.
On a balanced dataset this is already weak; on an imbalanced one it
tends to predict the majority class for most observations, which
yields high accuracy but very low recall for the minority, collapsing
macro F1.

\textbf{MNIST accuracy of 93\%} is impressive for a single unpruned
tree on raw pixels.  Digits 3 and 8 share curved strokes and are one of
the hardest MNIST binary pairs, yet 784 pixel features provide enough
discriminative signal for a deep tree.  Our custom tree achieves 92.9\%
vs.\ sklearn's 93.2\%, a difference of 0.25\%.

\textbf{Covertype correctness.} The custom tree matches sklearn
within 0.22\% (0.7802 vs. 0.7780), confirming the multiclass split-search
generalises correctly beyond binary problems. The stump collapses to 63.88\%
accuracy but only 0.1968 macro-F1--a single binary split cannot
meaningfully separate seven classes, foreshadowing the same
capacity limitation that drives AdaBoost's macro-F1 collapse in
Experiment 4.

\subsection{Experiment 2: AdaBoost Scaling}
\label{ssec:ada_scale}

Figure~\ref{fig:adaboost_scaling} shows test and train accuracy
as a function of the number of stumps $T$.

\begin{figure}[h]
\centering
\includegraphics[width=0.95\linewidth]{figures/adaboost_scaling_breast_cancer.png}
\caption{AdaBoost train/test accuracy vs.\ number of stumps on Breast Cancer.
Train accuracy quickly reaches 1.0; test accuracy plateaus around $T{=}30$.}
\label{fig:adaboost_scaling}
\end{figure}

\textbf{Observation.}  On Breast Cancer, test accuracy plateaus near
$0.97$ by $T{\approx}30$ stumps; additional estimators do not help.
On Adult Income (Fig.~\ref{fig:adaboost_adult}), convergence is
slower because the boundary is more complex, and the training
accuracy reaches 1.0 at $T{\approx}80$.

\begin{figure}[h]
\centering
\includegraphics[width=0.95\linewidth]{figures/adaboost_scaling_adult.png}
\caption{AdaBoost scaling on Adult Income.  Slower convergence due to
high-dimensional, imbalanced feature space.}
\label{fig:adaboost_adult}
\end{figure}

\textbf{Why does AdaBoost not overfit here?}  For exponential loss
and stumps, Schapire showed that the margin distribution keeps
improving even after training error reaches zero, explaining the
continued test accuracy improvement before the plateau.

\subsection{Experiment 3: Random Forest Scaling}
\label{ssec:rf_scale}

\begin{figure}[h]
\centering
\includegraphics[width=0.95\linewidth]{figures/rf_estimators_scaling_breast_cancer.png}
\caption{Random Forest accuracy vs.\ number of trees (Breast Cancer).
Test accuracy stabilises near $T{=}50$.}
\label{fig:rf_estimators}
\end{figure}

\begin{figure}[h]
\centering
\includegraphics[width=0.95\linewidth]{figures/rf_depth_scaling_breast_cancer.png}
\caption{Random Forest accuracy vs.\ max tree depth (Breast Cancer).
Accuracy peaks around \texttt{max\_depth}=10; deeper trees overfit.}
\label{fig:rf_depth}
\end{figure}

\textbf{Number of trees} (Fig.~\ref{fig:rf_estimators}).  Accuracy
is noisy for $T{<}20$ trees (high variance from limited averaging)
then smooths out above $T{=}50$.  Beyond $T{=}100$ there is
negligible gain-variance reduction saturates as the law of large
numbers kicks in.

\textbf{Tree depth} (Fig.~\ref{fig:rf_depth}).  Shallow trees
($\text{depth}\le 4$) underfit; deep trees ($\text{depth}\ge 15$)
memorise training data without proportional test improvement.  The
sweet spot is depth $\approx 10$ for the Breast Cancer dataset.

\textbf{Covertype requires substantially more capacity.} Unlike
Breast Cancer's early plateau ($\approx$depth 10), Covertype's test
accuracy is still rising at depth 20 ($0.526 \to 0.799$), and
estimator scaling plateaus later ($T\approx$40-60 vs.\ $T\approx$50
for Breast Cancer). This confirms that optimal Random Forest
capacity is problem-complexity-dependent, not a fixed default.
This contrasts with AdaBoost: stumps (depth~1) work for AdaBoost
because the ensemble aggregates thousands of complementary cuts.

\subsection{Experiment 4: Head-to-Head Comparison}
\label{ssec:h2h}

Figures~\ref{fig:h2h_bc}-\ref{fig:h2h_cov_f1} show box plots from
5-fold cross-validation at $T{=}100$ estimators.

\begin{figure}[h]
\centering
\includegraphics[width=0.95\linewidth]{figures/head_to_head_accuracy_breast_cancer.png}
\caption{5-fold CV accuracy on Breast Cancer.  Both ensembles
substantially outperform the single tree with lower variance.}
\label{fig:h2h_bc}
\end{figure}

\begin{figure}[h]
\centering
\includegraphics[width=0.95\linewidth]{figures/head_to_head_auc_roc_breast_cancer.png}
\caption{5-fold CV AUC-ROC on Breast Cancer.  Random Forest achieves
the highest median AUC-ROC with the smallest inter-quartile range.}
\label{fig:h2h_auc}
\end{figure}

\begin{figure}[h]
\centering
\includegraphics[width=0.95\linewidth]{figures/head_to_head_accuracy_mnist.png}
\caption{5-fold CV accuracy on MNIST 3-vs-8.  High dimensionality
benefits both ensemble methods equally.}
\label{fig:h2h_mnist}
\end{figure}

\begin{figure}[h]
\centering
\includegraphics[width=0.95\linewidth]{figures/head_to_head_accuracy_covertype.png}
\caption{5-fold CV accuracy on Covertype under 7-way imbalance.}
\label{fig:h2h_cov_acc}
\end{figure}

\begin{figure}[h]
\centering
\includegraphics[width=0.95\linewidth]{figures/head_to_head_macro_f1_covertype.png}
\caption{5-fold CV macro F1-score on Covertype, highlighting AdaBoost's F1 collapse.}
\label{fig:h2h_cov_f1}
\end{figure}

\textbf{Key findings.}  (1) Both ensembles reduce inter-quartile
range significantly compared with the single tree, confirming
variance reduction.  (2) On Breast Cancer, AdaBoost achieves the
highest mean accuracy ($0.9755 \pm 0.0186$) while Random Forest
is close behind ($0.9631 \pm 0.0203$); on AUC-ROC both reach
$>0.98$.  (3) On Adult Income, our custom Random Forest
($0.8623 \pm 0.0032$) marginally outperforms sklearn's RF
($0.8507 \pm 0.0023$), demonstrating that the custom implementation
is competitive.  (4) On MNIST with 784 raw-pixel
features, Random Forest dominates with $0.9808 \pm 0.0035$ accuracy
vs.\ AdaBoost's $0.9505 \pm 0.0032$-the random feature
sub-sampling effectively focuses each tree on the most
informative pixel regions.

\textbf{(5) On severely imbalanced Covertype, accuracy alone is
misleading.} AdaBoost reaches 0.6256 accuracy but only
\textbf{0.3186} macro-F1--a 0.31 gap indicating the model has
learned to ignore most of the seven classes in favour of the
dominant ones. Random Forest (sklearn) achieves both higher accuracy
(0.8501) \emph{and} a far smaller accuracy-F1 gap (0.7457, a 0.10
gap), showing bagging tolerates severe multiclass imbalance far
better than sequential re-weighting. This is the clearest evidence
in our study that boosting's aggressive bias reduction can actively
work against it when the class distribution itself is the dominant
source of difficulty.

\subsection{Experiment 5: Noise Robustness}
\label{ssec:noise}

We randomly flip a fraction $\eta$ of training labels.

\begin{figure}[h]
\centering
\includegraphics[width=0.95\linewidth]{figures/noise_robustness_breast_cancer.png}
\caption{Accuracy vs.\ label noise fraction $\eta$ on Breast Cancer.
Random Forest degrades gracefully; AdaBoost collapses at high noise.}
\label{fig:noise_bc}
\end{figure}

\begin{figure}[h]
\centering
\includegraphics[width=0.95\linewidth]{figures/noise_robustness_adult.png}
\caption{Noise robustness on Adult Income.  Random Forest retains
a higher absolute accuracy at every $\eta$, but AdaBoost degrades
\emph{less} in relative terms ($\Delta{=}0.0019$ vs.\ RF's
$\Delta{=}0.0056$ from $\eta{=}0$ to $0.20$)-the opposite of the
degradation-rate ordering seen on Breast Cancer.}
\label{fig:noise_adult}
\end{figure}

\textbf{Why does AdaBoost collapse under noise?}  The re-weighting
mechanism assigns exponentially higher weight to misclassified points.
When some of those points are mislabelled (genuinely hard boundary
samples that are now labelled incorrectly), the stump trained in the
next round attempts to classify what are in fact random-label points,
degrading generalisation.  At $\eta{=}0.20$, approximately one fifth
of all re-weighted points carry corrupted labels, causing the ensemble
to memorise noise.

\textbf{Why does Random Forest stay stable?}  Each tree in the
forest is trained on a separate bootstrap sample, so each individual
tree is only partially exposed to the corrupted labels.  When
predictions are averaged, the impact of any single corrupted tree
is diluted.  Furthermore, since Random Forest does not focus more
resources on hard samples the way AdaBoost does, it avoids the
amplification effect.  This amplification is most pronounced on
Breast Cancer, where the class boundary is clean; on Adult Income,
where the boundary is already noisy and complex, AdaBoost's
re-weighting degrades more gracefully than on cleaner data,
yielding a smaller accuracy drop than Random Forest at $\eta{=}0.20$.

\begin{figure}[h]
\centering
\includegraphics[width=0.95\linewidth]{figures/noise_robustness_covertype.png}
\caption{Noise robustness on Covertype under label corruption.}
\label{fig:noise_cov}
\end{figure}

\textbf{An anomaly on Covertype.} Unlike the other three datasets,
AdaBoost's accuracy on Covertype \emph{increases} under label
noise ($0.5927\to0.6832$ from $\eta{=}0$ to $\eta{=}0.20$; all runs
fit the full 100-estimator budget, ruling out early stopping as the
cause). Class-wise recall diagnostics confirm this is a noise-induced
regularisation effect under severe imbalance. At $\eta{=}0.00$,
AdaBoost is forced to focus on misclassified minority instances, achieving
moderate recall on them ($40.0\%$ on Class 5, $7.73\%$ on Class 6) at the
cost of distorting dominant-class boundaries (Class 1 recall is only $53.6\%$).
Under label noise ($\eta{=}0.05$), the rare classes become unlearnable, so
AdaBoost ``gives up'' on them (Class 5 recall drops to $1.1\%$, Class 6 to $0\%$)
and focuses exclusively on the dominant classes (Classes 0, 1, 2). This significantly
improves prediction quality on these classes (Class 1 recall climbs to $74.8\%$;
Class 2 recall climbs to $71.0\%$). Since the dominant classes constitute $91.5\%$
of the dataset, this trade-off yields a net $+8.5\%$ gain in overall test accuracy.

\subsection{Experiment 6: Bias-Variance Decomposition}
\label{ssec:biasvar}

Using $B{=}100$ bootstrap training sets, we estimate Bias$^2$ and
Variance following Breiman~\cite{breiman1996bagging}
(Table~\ref{tab:biasvar} and Fig.~\ref{fig:biasvar}).

\begin{table}[h]
\centering
\caption{Bias-Variance decomposition (Breast Cancer, $B{=}100$).}
\label{tab:biasvar}
\begin{tabular}{lcc}
\toprule
\textbf{Model} & \textbf{Bias$^2$} & \textbf{Variance}\\
\midrule
Single Tree   & 0.0408 & 0.0352\\
AdaBoost      & \textbf{0.0231} & 0.0131\\
Random Forest & 0.0317 & \textbf{0.0076}\\
\bottomrule
\end{tabular}
\end{table}

\begin{figure}[h]
\centering
\includegraphics[width=0.95\linewidth]{figures/bias_variance.png}
\caption{Bias$^2$ and Variance for three models on Breast Cancer.
AdaBoost wins on bias; Random Forest wins on variance.}
\label{fig:biasvar}
\end{figure}

\textbf{Interpretation.}
The single tree has the highest values on both dimensions: it overfits
some bootstrap samples (high variance) while still making systematic
boundary errors (high bias).

AdaBoost reduces Bias$^2$ by $43\%$ relative to the tree (from
0.0408 to 0.0231) by iteratively correcting systematic errors.  Its
Variance is moderate (0.0131) because each round's stump is
deterministic given the weights, but the overall ensemble is sensitive
to which samples are upweighted, introducing some cross-replicate
variation.

Random Forest achieves the lowest Variance (0.0076, a $78\%$
reduction from the single tree) through averaging of decorrelated
trees.  Its Bias$^2$ (0.0317) is intermediate: each tree is
unbiased (full-depth), but feature sub-sampling means individual
splits may not always choose the globally optimal feature, leaving
some systematic error.

\textbf{Total error.}  AdaBoost: $0.0231+0.0131=0.0362$.  Random
Forest: $0.0317+0.0076=0.0393$.  On this dataset AdaBoost wins
marginally on total decomposed error.

\subsection{Experiment 7: Unsupervised Analysis}
\label{ssec:unsup}

\begin{figure}[h]
\centering
\includegraphics[width=0.95\linewidth]{figures/scree_plot.png}
\caption{PCA scree plot on Breast Cancer.  90\% of variance is captured
by the first 7 principal components.}
\label{fig:scree}
\end{figure}

\begin{figure}[h]
\centering
\includegraphics[width=0.95\linewidth]{figures/elbow_method.png}
\caption{K-Means elbow plot.  Inertia decreases sharply up to $k{=}2$,
confirming the binary class structure.}
\label{fig:elbow}
\end{figure}

\begin{figure}[h]
\centering
\includegraphics[width=0.95\linewidth]{figures/pca_scatter_comparison.png}
\caption{2-D PCA projection coloured by true labels (left), K-Means
clusters (centre), and DBSCAN clusters (right).}
\label{fig:pcascatter}
\end{figure}

\textbf{PCA.}  The first 7 of 30 principal components account for
$\ge 90\%$ of the variance (Fig.~\ref{fig:scree}).  This moderate
effective dimensionality explains why even a single unpruned tree
achieves over 94\% accuracy-the majority of 30 features are
correlated cell-nucleus measurements, and the 7-PC projection
captures the essential discriminative structure.

\textbf{K-Means.}  The elbow at $k{=}2$ (Fig.~\ref{fig:elbow})
aligns with the binary classification task.  The K-Means clusters
recover the malignant/benign separation with an Adjusted Rand Index
(ARI) of $\mathbf{0.6536}$.  The moderate ARI reflects that the two classes
partially overlap in PCA space (Fig.~\ref{fig:pcascatter}), which
is the same overlap that challenges classifiers.

\textbf{DBSCAN.}  With $\varepsilon{=}3.5$ (chosen from the
k-distance plot, Fig.~\ref{fig:kdist}), DBSCAN achieves ARI of
$0.0653$, substantially lower than K-Means.  This low ARI indicates
that density-based clustering struggles on the Breast Cancer dataset:
the two classes are not separated by a density gap but rather blend
along a continuous manifold.  DBSCAN additionally labels
approximately $13.2\%$ of samples as noise-boundary-region points
corresponding to the clinically ambiguous cases that are hardest
to classify correctly.

\begin{figure}[h]
\centering
\includegraphics[width=0.95\linewidth]{figures/k_distance_plot.png}
\caption{K-distance plot for DBSCAN parameter selection.  The knee
around $\varepsilon{\approx}3.5$ guides our choice.}
\label{fig:kdist}
\end{figure}

\subsection{Bonus: SAMME vs.\ SAMME.R}
\label{ssec:sammer}

\begin{figure}[h]
\centering
\includegraphics[width=0.95\linewidth]{figures/sammer_comparison.png}
\caption{Calibration comparison: SAMME (discrete) vs.\ SAMME.R (real)
across datasets.  SAMME.R converges faster and produces better-calibrated
probabilities.}
\label{fig:sammer}
\end{figure}

SAMME.R uses probability estimates from each stump rather than hard
predictions, typically converging faster and yielding better-calibrated
probabilities.  Figure~\ref{fig:sammer} confirms that SAMME.R achieves
lower log-loss at the same number of estimators (0.1322 vs.\ 0.6648 on
Iris, at identical 90\% test accuracy for both variants).

\subsection{Bonus: Gradient Boosting vs.\ AdaBoost}
\label{ssec:gbm}

\begin{figure}[h]
\centering
\includegraphics[width=0.95\linewidth]{figures/gbm_vs_adaboost_breast_cancer.png}
\caption{AdaBoost vs.\ GBM test accuracy on Breast Cancer.  AdaBoost
converges to a higher plateau despite using only depth-1 stumps.}
\label{fig:gbm}
\end{figure}

Our custom Gradient Boosting implementation fits trees on negative
gradients of the log-loss.  AdaBoost outperforms it on all three
binary datasets (Fig.~\ref{fig:gbm}): Breast Cancer (0.9683 vs.\
0.9561), Adult Income (0.8310 vs.\ 0.8300), and MNIST (0.9450 vs.\
0.9270).  This is somewhat counter to the usual expectation that
deeper trees optimising log-loss should out-perform stumps
optimising exponential loss.  We attribute this to the fixed budget
used here (100 estimators, learning rate 0.1, depth 3): AdaBoost's
stumps are extremely low-capacity (low VC dimension), so its
adaptive re-weighting extracts most of the available signal within
100 rounds, while GBM's deeper, higher-capacity trees at this
learning rate and round count have not yet fully converged in
log-loss space.  A higher learning rate or more estimators would
likely narrow or reverse this gap.  Covertype is excluded from this
comparison, as our Gradient Boosting implementation targets binary
classification.

% -- Discussion ----------
\section{Discussion}
\label{sec:discussion}

\subsection{When Does Boosting Win?}
On \emph{clean, small, and well-structured} datasets (Breast Cancer),
AdaBoost matches or slightly exceeds Random Forest in total Bias$^2+$Variance
(0.0362 vs.\ 0.0393).  Its sequential focus on hard samples drives the
decision boundary into complex regions that naive averaging misses.

\subsection{When Does Bagging Win?}
Random Forest dominates in \emph{four} scenarios:
\begin{enumerate}
  \item \textbf{Label noise (on clean-boundary data).}  Averaging
        independent trees dilutes corrupted-label influence;
        AdaBoost amplifies it when the underlying boundary is
        otherwise clean, as on Breast Cancer.
  \item \textbf{High variance.}  On smaller datasets (or noisier
        features) the bootstrapped averaging provides stability
        that sequential re-weighting cannot.
  \item \textbf{High-dimensional feature spaces.}  Random feature
        sub-sampling at each node effectively performs implicit
        feature selection per tree, making Random Forest naturally
        robust to irrelevant features without pre-processing.
  \item \textbf{Severe multiclass imbalance.} On Covertype, Random
        Forest's accuracy-macro-F1 gap (0.10) is roughly a third of
        AdaBoost's (0.31), showing bagging degrades more gracefully than
        boosting when classes are extremely unevenly represented.
\end{enumerate}
Notably, Covertype's noise-robustness result (\S\ref{ssec:noise}) shows
that AdaBoost's noise sensitivity is not universal: under severe
imbalance, injected noise can \emph{help} AdaBoost by disrupting a
degenerate rare-class-fitting regime, the opposite of its behaviour
on clean data.  This qualifies scenario (1) above as specific to
otherwise-clean boundaries rather than a universal property of
boosting.

\subsection{Connection to Unsupervised Structure}
The PCA analysis shows that Breast Cancer's variance is spread across
7 principal components (not a tight 2D cluster), yet both ensembles
achieve near-perfect AUC-ROC ($\ge 0.98$).  The high K-Means ARI
($0.6536$) confirms that the malignant/benign boundary is largely
linearly separable in PC space.  By contrast, DBSCAN's low ARI
($0.0653$) and noise fraction of $13.2\%$ reveal that the
density structure is diffuse-there is no clean density gap between
classes.  The $13.2\%$ noise points are the boundary-region samples
that AdaBoost re-weights heavily in later rounds, which explains
both its lower bias and its vulnerability to label noise on this
particular dataset: it cares disproportionately about the hardest
cases.

\subsection{Implementation Correctness}
Our custom tree matches sklearn within 1\% absolute accuracy on all
four datasets (Table~\ref{tab:baseline}), confirming that the $O(N\log N)$
cumsum sweep correctly implements the CART criterion without
floating-point errors.

% -- Conclusion ----------
\section{Conclusion}
\label{sec:conclusion}

We implemented CART, AdaBoost (SAMME), and Random Forest from scratch
and evaluated them across seven controlled experiments.  The key
finding is that neither algorithm uniformly dominates: \emph{AdaBoost
reduces Bias$^2$ more aggressively; Random Forest reduces Variance
more dramatically and is significantly more robust to label corruption
on otherwise-clean data and to severe multiclass imbalance.}
Practitioners should prefer boosting when data is clean and structured,
and bagging when noise, class imbalance, or high dimensionality is a
concern.  A secondary finding is that noise robustness itself is
context-dependent: AdaBoost's sensitivity to label corruption, while
strong on clean data, can invert under severe class imbalance, where
injected noise disrupts a degenerate rare-class-fitting regime and
improves accuracy.

\textbf{Limitations.}  Experiments use fixed hyperparameters
(100 estimators, \texttt{max\_features=sqrt}).  A full grid search
might narrow the performance gap.  All experiments use the full or subsampled
datasets: Breast Cancer (569 samples), Adult Income (full 45,222
examples after dropping rows with missing values via \texttt{dropna}), MNIST
3-vs-8 (6,000 samples), and Covertype (30,000-sample random subsample,
preserving the full dataset's $\approx$0.5\% minority-class rarity).
Our Gradient Boosting bonus targets binary classification only and
was therefore not evaluated on Covertype.


% ===== BIBLIOGRAPHY ===================================================
\begin{thebibliography}{00}

\bibitem{hastie2009elements}
T.~Hastie, R.~Tibshirani, and J.~Friedman,
\textit{The Elements of Statistical Learning}, 2nd~ed.
Springer, 2009.

\bibitem{quinlan1986induction}
J.~R.~Quinlan,
``Induction of decision trees,''
\textit{Machine Learning}, vol.~1, no.~1, pp.~81-106, 1986.

\bibitem{freund1997decision}
Y.~Freund and R.~E.~Schapire,
``A decision-theoretic generalization of on-line learning and an
application to boosting,''
\textit{J. Comput. Syst. Sci.}, vol.~55, no.~1, pp.~119-139, 1997.

\bibitem{breiman1996bagging}
L.~Breiman,
``Bagging predictors,''
\textit{Machine Learning}, vol.~24, no.~2, pp.~123-140, 1996.

\bibitem{breiman2001random}
L.~Breiman,
``Random forests,''
\textit{Machine Learning}, vol.~45, no.~1, pp.~5-32, 2001.

\bibitem{pearson1901lines}
K.~Pearson,
``On lines and planes of closest fit to systems of points in space,''
\textit{Philosophical Magazine}, vol.~2, no.~11, pp.~559-572, 1901.

\bibitem{lloyd1982least}
S.~P.~Lloyd,
``Least squares quantization in PCM,''
\textit{IEEE Trans. Inf. Theory}, vol.~28, no.~2, pp.~129-137, 1982.

\bibitem{ester1996density}
M.~Ester, H.-P.~Kriegel, J.~Sander, and X.~Xu,
``A density-based algorithm for discovering clusters in large spatial
databases with noise,''
in \textit{Proc. 2nd Int. Conf. Knowledge Discovery and Data Mining
(KDD)}, 1996, pp.~226-231.

\end{thebibliography}

\end{document}